% Table for Section 5.2.2 (optional) — mean normalized margin of pair 1, 10 and 30 over all contexts of each error type.
% Generated by result_5_2.ipynb, Section 7 (numbers: data/fig_5_11_margin_by_pair.csv).
\begin{table}[!t]
    \centering
    \footnotesize
    \setlength{\tabcolsep}{5pt}
    \renewcommand{\arraystretch}{1.08}
    \captionsetup{skip=4pt}
    \begin{tabularx}{\linewidth}{@{}>{\raggedright\arraybackslash}Xlrrrr@{}}
        \toprule
        Student error type & Move & Contexts & $j = 1$ & $j = 10$ & $j = 30$ \\
        \midrule
        Numerical calculation error & Focus & 87 & 0.71 & 0.27 & 0.04 \\
        None of the above & Focus & 60 & 0.71 & 0.28 & 0.04 \\
        Reached the correct solution but proceeded further & Focus & 49 & 0.49 & 0.11 & 0.01 \\
        Missing or extra quantity & Telling & 167 & 0.42 & 0.12 & 0.01 \\
        Missing or incorrect factual knowledge & Telling & 97 & 0.37 & 0.11 & 0.01 \\
        Unit conversion error & Telling & 34 & 0.48 & 0.17 & 0.02 \\
        Misunderstanding of the question & Generic & 199 & 0.03 & 0.00 & 0.00 \\
        \bottomrule
    \end{tabularx}
    \caption{Mean normalized score margin $\Delta_{e,j}/\max_{m}\operatorname{PPMI}(m,e)$ of the first, tenth, and 30th pair of a context, over all contexts of each error type (Figure~\ref{fig:margin_by_pair}). Move: the move with the largest utility of the error type.}
    \label{tab:margin_by_pair}
\end{table}
